\documentclass[11pt,a4paper]{article}

% ---------- Paquets ----------
\usepackage[utf8]{inputenc}
\usepackage[T1]{fontenc}
\usepackage[french]{babel}
\usepackage{lmodern}
\usepackage[margin=2.5cm]{geometry}
\usepackage{amsmath,amssymb}
\usepackage{icomma}
\usepackage{graphicx}
\usepackage{float}
\usepackage{caption}
\usepackage{array}
\usepackage{xcolor}
\usepackage{xspace}
\usepackage{pgfplots}
\pgfplotsset{compat=1.18}
\usetikzlibrary{babel}
\usepackage[hidelinks]{hyperref}

\definecolor{orangeesir}{RGB}{230,120,40}
\pgfplotsset{/pgf/number format/.cd, use comma, 1000 sep={\,}}
\renewcommand{\arraystretch}{1.2}
\newcommand{\us}{\,\textmu s\xspace}

\begin{document}

% =====================================================================
%                          PAGE DE GARDE
% =====================================================================
\pagenumbering{gobble}
\begin{titlepage}
\begin{tikzpicture}[remember picture,overlay]
  \draw[line width=0.6pt]
    ([xshift=2cm,yshift=-2cm]current page.north west)
    rectangle
    ([xshift=-2cm,yshift=2cm]current page.south east);
\end{tikzpicture}

\centering
\vspace*{0.3cm}

\includegraphics[width=0.45\textwidth]{logo_univ_rennes.png}

\vspace{1cm}

\includegraphics[width=0.22\textwidth]{logo esir.png}

\vspace{3cm}

{\Large\bfseries \underline{Travaux Pratiques 2 :} Transformée de Fourier rapide\par}
\vspace{0.3cm}
{\large Partie 1 : FFT récursive et multiplication de polynômes\par}
\vspace{0.3cm}
Algorithmique et Complexité, ESIR 2 IOT

\vspace{3cm}

\begin{minipage}[t]{0.42\textwidth}
  \centering
  {\color{orangeesir}\bfseries PROFESSEUR}\\[0.3cm]
  \fbox{\parbox[c][2cm][c]{0.9\linewidth}{\centering M. Guillaume Le GULUDEC}}
\end{minipage}
\hfill
\begin{minipage}[t]{0.42\textwidth}
  \centering
  {\color{orangeesir}\bfseries ÉTUDIANTS}\\[0.3cm]
  \fbox{\parbox[c][2cm][c]{0.9\linewidth}{\centering
    COULIBALY Yaplé Assétou\\[0.2cm]
    KOUAKOU Yah Mardochée}}
\end{minipage}

\vfill

\fbox{\parbox{0.6\textwidth}{\centering\bfseries\itshape ANNÉE ACADÉMIQUE : 2026-2027}}

\vspace{0.5cm}
\end{titlepage}

\pagenumbering{roman}
\tableofcontents
\newpage
\pagenumbering{arabic}

% =====================================================================

% =====================================================================
\section{Introduction}
% =====================================================================

Dans cette partie, on a programmé la FFT récursive et la TFD inverse, puis la
multiplication de polynômes par FFT. On a ensuite mesuré les temps de calcul de la
multiplication par FFT et de la multiplication par les coefficients, pour vérifier
qu'on retrouve bien $\Theta(n \log n)$ et $\Theta(n^2)$.

% =====================================================================
\section{FFT récursive et TFD inverse}
% =====================================================================

\subsection{Ce qu'on a fait}

La fonction de combinaison applique les formules de l'énoncé : à partir des TFD des
coefficients pairs et impairs (taille $n/2$), elle calcule les $n$ valeurs
$A(\omega_n^k)$ et $A(\omega_n^{k+n/2})$, avec le même produit
$\omega_n^k A_{\text{impair}}(\omega_{n/2}^k)$ pour les deux. La FFT sépare le vecteur
en indices pairs et impairs, s'appelle sur chaque moitié et combine. On s'arrête à
$n = 1$, où la TFD est le vecteur lui-même ; $n$ doit donc être une puissance de 2.
Pour l'inverse, on a utilisé la formule de l'énoncé : conjuguer, appliquer la FFT,
conjuguer et diviser par $n$.

\subsection{Complexité}

Chaque appel fait deux appels sur $n/2$ et une combinaison en $\Theta(n)$ :
$T(n) = 2\,T(n/2) + \Theta(n)$, donc $T(n) = \Theta(n \log n)$ ($\log_2 n$ niveaux
qui coûtent chacun $\Theta(n)$). L'inverse ajoute deux conjugaisons et une division
en $\Theta(n)$, elle reste en $\Theta(n \log n)$. Notre version crée deux tableaux et
calcule un cosinus et un sinus à chaque appel : ça ne change pas la complexité, mais
ça augmente la constante, ce qu'on retrouve dans les mesures.

\subsection{Tests}

Sur l'exemple du TD, $a = (1, 2, 3, 4)$, on attendait
$A(1) = 10$, $A(i) = -2 - 2i$, $A(-1) = -2$, $A(-i) = -2 + 2i$. Notre FFT renvoie
exactement $[\,10,\ -2-2i,\ -2,\ -2+2i\,]$, et l'inverse appliquée à ce résultat
redonne $[\,1,\ 2,\ 3,\ 4\,]$.

% =====================================================================
\section{Multiplication de polynômes}
% =====================================================================

On a suivi la méthode du TD : FFT de $A$ et de $B$, produit terme à terme, FFT
inverse, soit $\Theta(n \log n)$ au total. Le produit est de degré $2n-2$ : il faut
au moins $2n-1$ points, donc on complète les deux vecteurs avec des zéros jusqu'à la
taille $2n$ (qui reste une puissance de 2). On a testé les deux exemples de l'énoncé
et $A_3 \times B_3$ avec $A_3 = 1 + 2X + 3X^2 + 4X^3$ et $B_3 = -3 + 2X - 5X^2$, dont
le produit attendu est $-3 - 4X - 10X^2 - 16X^3 - 7X^4 - 20X^5$.

\begin{table}[H]
\centering
\caption{Tests de la multiplication par FFT.}
\label{tab:mult}
\begin{tabular}{|l|l|l|}
\hline
\textbf{Produit} & \textbf{Résultat attendu} & \textbf{Coefficients obtenus} \\
\hline
$2 \times (-3)$ & $-6$ & $[\,-6,\ 0\,]$ \\
\hline
$(2 + X)(-3 + 2X)$ & $-6 + X + 2X^2$ & $[\,-6,\ 1,\ 2,\ 0\,]$ \\
\hline
$A_3 \times B_3$ & voir ci-dessus &
  $[\,-3,\ -4,\ -10,\ -16,\ -7,\ -20,\ 0,\ 0\,]$ \\
\hline
\end{tabular}
\end{table}

Les résultats sont justes, et la méthode par les coefficients donne la même chose
pour $A_3 \times B_3$. Quelques remarques :
\begin{itemize}
  \item on obtient $2n$ valeurs pour $2n-1$ coefficients : la dernière vient du
        doublement de taille et vaut toujours 0 ;
  \item notre fonction suppose que les deux polynômes ont le même nombre de
        coefficients, on a donc écrit $B_3 = (-3, 2, -5, 0)$ ;
  \item les calculs sont faits en flottants, les coefficients sont justes à des
        erreurs d'arrondi près, trop petites pour apparaître à l'affichage.
\end{itemize}

% =====================================================================
\section{Comparaison des temps de calcul}
% =====================================================================

\subsection{Démarche}

On a mesuré le temps de multiplication de deux polynômes aléatoires avec les deux
méthodes, pour $n = 2^6$ à $2^{15}$, en microsecondes. Avant la série, on fait
20 multiplications non mesurées pour que la machine virtuelle Java ait déjà compilé
le code, et pour chaque taille on garde le minimum de 5 exécutions.

Les modèles viennent de la complexité théorique :
\[
  T_{\text{coeff}}(n) = a\,n^2, \qquad T_{\text{FFT}}(n) = b\,n \ln n .
\]
Pour estimer $a$ et $b$ avec toutes les mesures et pas un seul point, on a calculé
$T(n)/n^2$ et $T(n)/(n \ln n)$ pour chaque taille. On a aussi calculé de combien le
temps est multiplié quand $n$ double : on attend 4 pour $n^2$ et
$2 \ln n / \ln(n/2)$, un peu plus de 2, pour $n \ln n$.

\subsection{Méthode par les coefficients}

\begin{table}[H]
\centering
\caption{Méthode par les coefficients (minimum de 5 exécutions).}
\label{tab:coeff}
\begin{tabular}{|r|r|r|r|}
\hline
$n$ & $T(n)$ (\textmu s) & $T(n)/n^2$ ($10^{-3}$\,\textmu s) & $T(n)/T(n/2)$ \\
\hline
64      & 17           & 4,15 & \\
128     & 70           & 4,27 & 4,12 \\
256     & 289          & 4,41 & 4,13 \\
512     & 1\,180       & 4,50 & 4,08 \\
1\,024  & 4\,822       & 4,60 & 4,09 \\
2\,048  & 18\,733      & 4,47 & 3,88 \\
4\,096  & 71\,680      & 4,27 & 3,83 \\
8\,192  & 276\,822     & 4,12 & 3,86 \\
16\,384 & 976\,377     & 3,64 & 3,53 \\
32\,768 & 3\,913\,325  & 3,64 & 4,01 \\
\hline
\end{tabular}
\end{table}

Le temps est multiplié par 3,5 à 4,1 quand $n$ double (3,95 en moyenne) : c'est bien
une croissance en $n^2$. $T(n)/n^2$ reste entre $4{,}1$ et $4{,}6 \cdot 10^{-3}$\us
jusqu'à $n = 8\,192$, puis descend à $3{,}6\cdot 10^{-3}$\us. Avec
$a = 0{,}0044$\us (valeur des petites tailles), le modèle passait 20\,\% au-dessus des
deux derniers points, on a donc pris
\[
  T_{\text{coeff}}(n) \approx 0{,}004 \cdot n^2 \ \text{\textmu s},
\]
qui reste à moins de 13\,\% de toutes les mesures (figure~\ref{fig:coeff}). On
revient sur la baisse du rapport dans la partie~\ref{sec:ecarts}.

\begin{figure}[H]
\centering
\begin{tikzpicture}
\begin{axis}[
  width=0.75\linewidth, height=6.5cm,
  xlabel={Taille $n$}, ylabel={Temps (\textmu s)},
  xmin=0, xmax=34000, ymin=0,
  xtick={0,8192,16384,24576,32768}, scaled x ticks=false,
  grid=major, legend pos=north west,
]
\addplot[only marks, mark=*, blue] coordinates {
  (64,17) (128,70) (256,289) (512,1180) (1024,4822) (2048,18733)
  (4096,71680) (8192,276822) (16384,976377) (32768,3913325)
};
\addlegendentry{Mesures}
\addplot[green!60!black, thick, domain=0:33000, samples=100] {0.004*x^2};
\addlegendentry{$T(n) = 0{,}004\, n^2$}
\end{axis}
\end{tikzpicture}
\caption{Méthode par les coefficients : mesures et modèle quadratique.}
\label{fig:coeff}
\end{figure}

\subsection{Méthode par FFT}

\begin{table}[H]
\centering
\caption{Méthode par FFT (minimum de 5 exécutions).}
\label{tab:fft}
\begin{tabular}{|r|r|r|r|r|}
\hline
$n$ & $T(n)$ (\textmu s) & $T(n)/(n \ln n)$ (\textmu s) & $T(n)/T(n/2)$ & attendu \\
\hline
64      & 107     & 0,402 &      &      \\
128     & 234     & 0,377 & 2,19 & 2,33 \\
256     & 512     & 0,361 & 2,19 & 2,29 \\
512     & 1\,861  & 0,583 & 3,63 & 2,25 \\
1\,024  & 2\,926  & 0,412 & 1,57 & 2,22 \\
2\,048  & 5\,276  & 0,338 & 1,80 & 2,20 \\
4\,096  & 6\,767  & 0,199 & 1,28 & 2,18 \\
8\,192  & 13\,590 & 0,184 & 2,01 & 2,17 \\
16\,384 & 28\,445 & 0,179 & 2,09 & 2,15 \\
32\,768 & 60\,935 & 0,179 & 2,14 & 2,14 \\
\hline
\end{tabular}
\end{table}

Jusqu'à $n = 4\,096$, les mesures sont irrégulières : le rapport tourne autour de 0,4
avec un saut à 0,58, et le temps est multiplié par 3,63, puis 1,57, puis 1,28. Un
facteur 1,28 quand $n$ double est impossible pour un algorithme en $n \ln n$ si le code
ne change pas, donc ces mesures sont faussées (partie~\ref{sec:ecarts}). À partir de
$n = 8\,192$, le rapport est stable (0,184 ; 0,179 ; 0,179) et les facteurs sont
proches de ceux attendus. On a pris $b$ dans cette zone :
\[
  T_{\text{FFT}}(n) \approx 0{,}18 \cdot n \ln n \ \text{\textmu s},
\]
à moins de 3\,\% des mesures pour $n \geq 8\,192$ (figure~\ref{fig:fft}).

\begin{figure}[H]
\centering
\begin{tikzpicture}
\begin{axis}[
  width=0.75\linewidth, height=6.5cm,
  xlabel={Taille $n$}, ylabel={Temps (\textmu s)},
  xmin=0, xmax=34000, ymin=0,
  xtick={0,8192,16384,24576,32768}, scaled x ticks=false,
  grid=major, legend pos=north west,
]
\addplot[only marks, mark=*, blue] coordinates {
  (64,107) (128,234) (256,512) (512,1861) (1024,2926) (2048,5276)
  (4096,6767) (8192,13590) (16384,28445) (32768,60935)
};
\addlegendentry{Mesures}
\addplot[orange, thick, domain=64:33000, samples=100] {0.18*x*ln(x)};
\addlegendentry{$T(n) = 0{,}18\, n \ln n$}
\end{axis}
\end{tikzpicture}
\caption{Méthode par FFT : mesures et modèle en $n \ln n$.}
\label{fig:fft}
\end{figure}

\subsection{Points qui s'écartent des modèles}
\label{sec:ecarts}

\begin{itemize}
  \item \textbf{FFT, petites tailles.} Les mesures jusqu'à $n = 4\,096$ sont au-dessus
        du modèle et irrégulières. On pense que c'est la compilation à la volée de
        Java : notre échauffement de 20 multiplications ne suffit sans doute pas pour
        les parties du code qui ne tournent qu'une fois par multiplication
        (conjugaison, produit terme à terme, recopie avec les zéros), et elles sont
        encore lentes au début de la série. On ne s'appuie donc que sur les grandes
        tailles pour estimer $b$.
  \item \textbf{Coefficients, changement de niveau.} $T(n)/n^2$ passe d'environ
        $4{,}4$ à $3{,}6 \cdot 10^{-3}$\us entre $n = 8\,192$ et $n = 16\,384$, alors
        que l'algorithme est le même. Une explication possible
        est un changement de fréquence du processeur (ordinateur portable) pendant la
        série. L'écart reste petit devant la croissance du temps, multiplié par plus de
        200\,000 entre $n = 64$ et $n = 32\,768$.
\end{itemize}

\subsection{Comparaison des deux méthodes}

\begin{figure}[H]
\centering
\begin{tikzpicture}
\begin{loglogaxis}[
  width=0.75\linewidth, height=6.5cm,
  xlabel={Taille $n$}, ylabel={Temps (\textmu s)},
  log basis x=2,
  grid=major, legend pos=north west,
]
\addplot[blue, mark=*] coordinates {
  (64,17) (128,70) (256,289) (512,1180) (1024,4822) (2048,18733)
  (4096,71680) (8192,276822) (16384,976377) (32768,3913325)
};
\addlegendentry{Par les coefficients}
\addplot[orange, mark=square*] coordinates {
  (64,107) (128,234) (256,512) (512,1861) (1024,2926) (2048,5276)
  (4096,6767) (8192,13590) (16384,28445) (32768,60935)
};
\addlegendentry{Par FFT}
\end{loglogaxis}
\end{tikzpicture}
\caption{Comparaison des deux méthodes (échelles logarithmiques : $n^2$ donne une
droite de pente 2, $n \ln n$ une pente un peu supérieure à 1).}
\label{fig:comparaison}
\end{figure}

Pour les petites tailles, la FFT est plus lente : 107\us contre 17\us à $n = 64$. Les
courbes se croisent entre $n = 512$ (1\,180\us contre 1\,861\us) et $n = 1\,024$
(4\,822\us contre 2\,926\us). C'est la constante : $b/a \approx 45$, à cause de la
récursion, des tableaux créés et des cosinus et sinus à chaque appel.

Avec les modèles, le croisement est là où $a\,n^2 = b\,n \ln n$, soit
$n / \ln n = b/a \approx 45$, ce qui donne $n \approx 250$. Nos mesures donnent un
croisement plus tard, entre 512 et 1\,024, parce que les temps de la FFT aux petites
tailles sont au-dessus du modèle (partie~\ref{sec:ecarts}).

Ensuite l'écart grandit comme $n / \ln n$ : à $n = 32\,768$, la FFT est 64 fois plus
rapide (61\,ms contre 3,9\,s).

\subsection{Extrapolation}

En prolongeant nos deux modèles jusqu'à $n = 2^{20}$ (environ un million de
coefficients) :
\[
  T_{\text{coeff}}(2^{20}) \approx 0{,}004 \times (2^{20})^2\ \text{\textmu s} \approx 73\ \text{min},
  \qquad
  T_{\text{FFT}}(2^{20}) \approx 0{,}18 \times 2^{20} \times \ln(2^{20})\ \text{\textmu s} \approx 2{,}6\ \text{s}.
\]
La FFT serait environ 1\,700 fois plus rapide. Ce sont des extrapolations de nos
modèles, pas des temps théoriques, et elles supposent
que $a$ et $b$ ne changent pas jusqu'à 32 fois notre plus grande taille (par exemple, les données ne tiennent plus dans le
cache du processeur) : seul l'ordre de grandeur compte.

% =====================================================================
\section{Conclusion}
% =====================================================================

Notre FFT et son inverse donnent les bons résultats, et la multiplication par FFT
donne les mêmes coefficients que la méthode classique. Les mesures confirment les
complexités : quand $n$ double, le temps est multiplié par environ 4 pour la méthode
par les coefficients et par un peu plus de 2 pour la FFT. Les points qui
s'écartent des modèles semblent venir de l'exécution du programme (compilation de
Java, fréquence du processeur) et pas des algorithmes. Enfin, à cause de sa constante 45 fois plus grande, la FFT
n'est plus rapide qu'à partir de quelques centaines de coefficients ; après, le gain
grandit comme $n / \ln n$.

\end{document}
